% c01b 原书 p.5–9（PDF p.012–016）
设 $A_i$（$i\in I$）是 $\sC$ 中一族对象。$A_i$（$i\in I$）的\emph{直积}\idx{zhiji}{直积}是一个元组 $(\prod_{i\in I}A_i,p_i)$，其中 $\prod_{i\in I}A_i$ 是 $\sC$ 中的对象，$p_i\colon\prod_{i\in I}A_i\to A_i$（$i\in I$）是一族态射，称为\emph{投影}\idx{touying}{投影}，使得以下泛性质成立：对任意对象 $C$ 与任意一族态射 $f_i\colon C\to A_i$（$i\in I$），存在唯一一个态射 $f\colon C\to\prod_{i\in I}A_i$，使 $p_if=f_i$ 对所有 $i$ 成立。若 $A_i$（$i\in I$）的直积存在，则它由唯一同构唯一确定，确切地说：若 $(\prod'_{i\in I}A_i,p'_i)$ 是另一个直积，则存在唯一的同构 $\phi\colon\prod_{i\in I}A_i\xrightarrow{\sim}\prod'_{i\in I}A_i$，使 $p'_i\phi=p_i$ 对所有 $i$ 成立。

$A_i$（$i\in I$）的\emph{直和}\idx{zhihe}{直和}用与上面类似的方式定义，只需把箭头的方向反过来。更确切地说，$A_i$（$i\in I$）的直和是一个元组 $(\bigoplus_{i\in I}A_i,k_i)$，其中 $\bigoplus_{i\in I}A_i$ 是 $\sC$ 中的对象，$k_i\colon A_i\to\bigoplus_{i\in I}A_i$（$i\in I$）是一族态射，使得以下泛性质成立：对任意对象 $C$ 与任意一族态射 $f_i\colon A_i\to C$（$i\in I$），存在唯一一个态射 $f\colon\bigoplus_{i\in I}A_i\to C$，使 $fk_i=f_i$ 对所有 $i$ 成立。若 $A_i$（$i\in I$）的直和存在，则它由唯一同构唯一确定。

设 $(I,\le)$ 是有向集。\emph{有向系统} $(A_i,\phi_{ij})_{i\in I}$ 由一族对象 $A_i$（$i\in I$）与态射 $\phi_{ij}\colon A_i\to A_j$（对满足 $i\le j$ 的下标 $i,j$）组成，使得 $\phi_{ii}=\id_{A_i}$ 对所有 $i$ 成立，且 $\phi_{jk}\phi_{ij}=\phi_{ik}$ 只要 $i\le j\le k$。有向系统 $(A_i,\phi_{ij})$ 的\emph{有向极限}是一个元组 $(\dlim_i A_i,\phi_i)$，其中 $\dlim_i A_i$ 是 $\sC$ 中的对象，$\phi_i\colon A_i\to\dlim_i A_i$（$i\in I$）是态射，满足 $\phi_j\phi_{ij}=\phi_i$ 只要 $i\le j$，使得以下泛性质成立：对任意对象 $C$ 与任意一族态射 $\psi_i\colon A_i\to C$（$i\in I$），若 $\psi_j\phi_{ij}=\psi_i$（$i\le j$），则存在唯一一个态射 $\psi\colon\dlim_i A_i\to C$，使 $\psi\phi_i=\psi_i$ 对所有 $i$ 成立。若有向极限存在，则它由唯一同构唯一确定。设 $(A'_i,\phi'_{ij})_{i\in I}$ 是另一个有向系统。从 $(A_i,\phi_{ij})$ 到 $(A'_i,\phi'_{ij})$ 的态射是一族态射 $u_i\colon A_i\to A'_i$（$i\in I$），使得对任意 $i\le j$，下图交换：
\begin{center}
\begin{tikzcd}
A_i \arrow[r,"u_i"] \arrow[d,"\phi_{ij}"'] & A'_i \arrow[d,"\phi'_{ij}"] \\
A_j \arrow[r,"u_j"'] & A'_j.
\end{tikzcd}
\end{center}
它诱导态射 $\dlim_i u_i\colon\dlim_i A_i\to\dlim_i A'_i$，使对所有 $i$，下图交换：
\begin{center}
\begin{tikzcd}
A_i \arrow[r,"u_i"] \arrow[d,"\phi_i"'] & A'_i \arrow[d,"\phi'_i"] \\
\dlim_i A_i \arrow[r,"\dlim_i u_i"'] & \dlim_i A'_i.
\end{tikzcd}
\end{center}

\emph{逆向系统} $(A_i,\phi_{ji})_{i\in I}$ 由一族对象 $A_i$（$i\in I$）与态射 $\phi_{ji}\colon A_j\to A_i$（对满足 $i\le j$ 的下标 $i,j$）组成，使得 $\phi_{ii}=\id_{A_i}$ 对任意 $i$ 成立，且 $\phi_{ji}\phi_{kj}=\phi_{ki}$ 只要 $i\le j\le k$。逆向系统 $(A_i,\phi_{ji})$ 的\emph{逆向极限}是一个元组 $(\plim_i A_i,\phi_i)$，其中 $\plim_i A_i$ 是 $\sC$ 中的对象，$\phi_i\colon\plim_i A_i\to A_i$（$i\in I$）是态射，满足 $\phi_{ji}\phi_j=\phi_i$ 只要 $i\le j$，使得以下泛性质成立：对任意对象 $C$ 与任意一族态射 $\psi_i\colon C\to A_i$（$i\in I$），若 $\phi_{ji}\psi_j=\psi_i$（$i\le j$），则存在唯一一个态射 $\psi\colon C\to\plim_i A_i$，使 $\phi_i\psi=\psi_i$ 对所有 $i$ 成立。从逆向系统 $(A_i,\phi_{ji})_{i\in I}$ 到逆向系统 $(A'_i,\phi'_{ji})_{i\in I}$ 的态射是一族态射 $u_i\colon A_i\to A'_i$（$i\in I$），使得对任意 $i\le j$，下图交换：
\begin{center}
\begin{tikzcd}
A_j \arrow[r,"u_j"] \arrow[d,"\phi_{ji}"'] & A'_j \arrow[d,"\phi'_{ji}"] \\
A_i \arrow[r,"u_i"'] & A'_i.
\end{tikzcd}
\end{center}
它诱导态射 $\plim_i u_i\colon\plim_i A_i\to\plim_i A'_i$。若 $(A_i,\phi_{ji})_{i\in I}$ 是集合的逆向系统，则 $\plim_i A_i$ 是由满足 $\phi_{ji}(x_j)=x_i$（$i\le j$）的元素 $(x_i)_{i\in I}\in\prod_i A_i$ 构成的 $\prod_i A_i$ 的子集，而 $\phi_i\colon\plim_i A_i\to A_i$ 是投影 $(x_i)_{i\in I}\mapsto x_i$。
\begin{tnote}
注意有向极限与逆向极限对集合的构造并不对称：有向极限是 $\coprod_{i\in I}A_i$ 的商（把同一元素在系统下的像等同起来），逆向极限则是 $\prod_{i\in I}A_i$ 的子集（挑出相容的元素族）。看到符号 $\dlim$ 与 $\plim$ 时先分清是哪一个，否则很容易把商与子集弄反。[依据：原书 p.2–3（有向极限的构造）与 p.6 第 1 段（逆向极限的构造）]
\end{tnote}

范畴 $\sC$ 称为\emph{加性范畴}\idx{jiaxingfanchou}{加性范畴}，若对 $\sC$ 中任意对象 $A,B,C$，$A$ 与 $B$ 的直积存在，$\Hom(A,B)$ 是交换群，且映射
\[
\Hom(A,B)\times\Hom(B,C)\to\Hom(A,C),\quad (f,g)\mapsto gf
\]
是双线性的。我们把 $0\in\Hom(A,B)$ 称为\emph{零态射}\idx{lingshe}{零态射}。

\begin{proposition}\label{it:1.1.2}
设 $\sC$ 是加性范畴，$A$ 与 $B$ 是 $\sC$ 中两个对象。
\begin{enumerate}[label=(\roman*),leftmargin=2.8em]
  \item 设 $p_1\colon A\times B\to A$ 与 $p_2\colon A\times B\to B$ 是投影。定义 $k_1\colon A\to A\times B$ 为满足 $p_1k_1=\id_A$ 与 $p_2k_1=0$ 的唯一态射，定义 $k_2\colon B\to A\times B$ 为满足 $p_1k_2=0$ 与 $p_2k_2=\id_B$ 的唯一态射。则 $k_1p_1+k_2p_2=\id_{A\times B}$。
  \item 设有一个对象 $P$ 与态射 $p_1\colon P\to A$，$p_2\colon P\to B$，$k_1\colon A\to P$，$k_2\colon B\to P$，使得
  \[
  p_1k_1=\id_A,\qquad p_2k_2=\id_B,\qquad k_1p_1+k_2p_2=\id_P.
  \]
  则 $P$ 连同态射 $p_1\colon P\to A$ 与 $p_2\colon P\to B$ 是 $A$ 与 $B$ 的直积，而 $P$ 连同态射 $k_1\colon A\to P$ 与 $k_2\colon B\to P$ 是 $A$ 与 $B$ 的直和。
\end{enumerate}
\end{proposition}

\begin{proof}
(i) 可以验证
\[
p_1(k_1p_1+k_2p_2)=p_1\id_{A\times B},\qquad
p_2(k_1p_1+k_2p_2)=p_2\id_{A\times B}.
\]
由直积的泛性质，得 $k_1p_1+k_2p_2=\id_{A\times B}$。

(ii) 我们有
\begin{align*}
p_1k_2&=p_1\id_Pk_2=p_1(k_1p_1+k_2p_2)k_2\\
&=(p_1k_1)(p_1k_2)+(p_1k_2)(p_2k_2)=p_1k_2+p_1k_2\\
&=2p_1k_2.
\end{align*}
故 $p_1k_2=0$。类似地 $p_2k_1=0$。

我们证明 $(P,k_1,k_2)$ 是 $A$ 与 $B$ 的直和，而把 $(P,p_1,p_2)$ 是 $A$ 与 $B$ 的直积的证明留给读者。给定对象 $C$ 与态射 $f_1\colon A\to C$、$f_2\colon B\to C$，定义 $f=f_1p_1+f_2p_2$。可以验证 $fk_1=f_1$ 且 $fk_2=f_2$。若 $f'\colon P\to C$ 是使 $f'k_1=f_1$ 且 $f'k_2=f_2$ 的态射，则
\begin{align*}
f'&=f'\id_P=f'(k_1p_1+k_2p_2)\\
&=(f'k_1)p_1+(f'k_2)p_2=f_1p_1+f_2p_2.
\end{align*}
这说明 $(P,k_1,k_2)$ 具有所需的泛性质。
\end{proof}
\begin{tnote}
命题 \ref{it:1.1.2} 是加性范畴与一般范畴的分水岭：它说明在加性范畴中，有限直积与有限直和作为对象是同一件东西（区别只在配上的投影或嵌入），所以后文对有限情形只写 $\oplus$，不再区分 $\prod$ 与 $\bigoplus$。注意结论 (i) 里的 $k_1p_1+k_2p_2$ 用到 $\Hom(A\times B,A\times B)$ 中的加法，这正是「加性」假设的用处。[依据：原书 p.6 命题 1.1.2 与 p.7 的证明]
\end{tnote}

设 $\sC$ 是加性范畴，$f\colon A\to B$ 是 $\sC$ 中的态射。若复合 $K\to A\to B$ 为 $0$，且对任意使复合 $K'\to A\to B$ 为 $0$ 的态射 $K'\to A$，存在唯一一个态射 $K'\to K$ 使下图交换，则称态射 $K\to A$ 是 $f$ 的\emph{核}\idx{he}{核}：
\begin{center}
\begin{tikzcd}
K' \arrow[d] \arrow[dr] & \\
K \arrow[r] & A
\end{tikzcd}
\end{center}
注意 $f$ 的核 $K\to A$ 必是单态射。我们常把 $K$ 记为 $\ker f$，并称它为 $f$ 的核。类似地，定义 $f$ 的\emph{余核}\idx{yuhe}{余核}为态射 $B\to C$，使得复合 $A\to B\to C$ 为 $0$，且对任意使复合 $A\to B\to C'$ 为 $0$ 的态射 $B\to C'$，存在唯一一个态射 $C\to C'$ 使下图交换：
\begin{center}
\begin{tikzcd}
B \arrow[r] \arrow[dr] & C \arrow[d] \\
& C'
\end{tikzcd}
\end{center}
$f$ 的余核 $B\to C$ 必是满态射。我们常把 $C$ 记为 $\coker f$，并称它为 $f$ 的余核。我们定义 $f$ 的\emph{像}\idx{xiang}{像}为 $f$ 的余核的核，定义 $f$ 的\emph{余像}\idx{yuxiang}{余像}为 $f$ 的核的余核。

存在从余像到像的典范态射 $\operatorname{coim}f\to\operatorname{im}f$，使下图交换：
\begin{center}
\begin{tikzcd}
A \arrow[r] \arrow[d] & B \\
\operatorname{coim}f \arrow[r] & \operatorname{im}f \arrow[u].
\end{tikzcd}
\end{center}
例如，当 $f\colon A\to B$ 是交换群范畴中的态射时，
\begin{align*}
\ker f&=\{a\in A\mid f(a)=0\},\\
\operatorname{im}f&=\{b\in B\mid \text{存在 } a\in A \text{ 使 } b=f(a)\},\\
\coker f&=B/\operatorname{im}f,\\
\operatorname{coim}f&=A/\ker f,
\end{align*}
而从余像到像的典范态射是典范同态 $A/\ker f\to\operatorname{im}f$（它是同构）。

设 $\sC$ 是加性范畴。$\sC$ 中的\emph{零对象}\idx{lingduixiang}{零对象} $0$ 是满足 $\Hom(0,0)=\{0\}$ 的对象。这等价于说 $0$ 的恒等态射等于零态射。对 $\sC$ 中任意对象 $X$，有 $\Hom(X,0)=\{0\}$ 与 $\Hom(0,X)=\{0\}$。$\sC$ 中的零对象彼此同构。

\emph{阿贝尔范畴}\idx{abeierfanchou}{阿贝尔范畴} $\sC$ 是具有零对象的加性范畴，使得对 $\sC$ 中每个态射 $f$，$f$ 的核与余核存在（从而 $f$ 的像与余像存在），且典范态射 $\operatorname{coim}f\to\operatorname{im}f$ 是同构。

在阿贝尔范畴中，双射态射是同构。事实上，若 $f\colon A\to B$ 是单射，则 $f$ 的核是 $0\to A$，$f$ 的余像是 $\id_A\colon A\to A$。若 $f\colon A\to B$ 是满射，则 $f$ 的余核是 $B\to 0$，$f$ 的像是 $\id_B\colon B\to B$。若 $f\colon A\to B$ 是双射，则典范态射 $\operatorname{coim}f\to\operatorname{im}f$ 恰是 $f\colon A\to B$。因 $\operatorname{coim}f\to\operatorname{im}f$ 是同构，$f\colon A\to B$ 是同构。
\begin{tnote}
「$\operatorname{coim}f\to\operatorname{im}f$ 是同构」这一条不能从其他公理推出，必须作为阿贝尔范畴的公理单独要求：它正是使「双射态射即同构」成立的条件。与此对照，$\cO_X$-模范畴与层范畴满足它（命题 \ref{it:1.1.6}），而一般加性范畴未必满足（见本节习题 \ref{ex:1.1.1}）。[依据：原书 p.8 第 2–4 段；反例见 p.17 习题 1]
\end{tnote}

设 $u\colon A\to B$ 是阿贝尔范畴中的单态射。此时称 $A$ 是 $B$ 的\emph{子对象}\idx{ziduixiang}{子对象}。设 $B\to C$ 是 $u$ 的余核。我们称 $C$ 为 $B$ 对 $A$ 的\emph{商}\idx{shang}{商}，记作 $B/A$。

在阿贝尔范畴中，态射序列
\[
A\xrightarrow{u}B\xrightarrow{v}C
\]
称为\emph{正合}，若 $vu=0$ 且典范态射 $\operatorname{coim}u\to\ker v$ 是同构。形如
\[
0\to A\xrightarrow{u}B\xrightarrow{v}C\to 0
\]
的正合序列称为\emph{短正合序列}\idx{duanzhenghexulie}{短正合序列}。若该短正合序列同构于
\[
0\to A\to A\oplus C\to C\to 0,
\]
则称它是\emph{分裂}的，其中 $A\to A\oplus C$ 与 $A\oplus C=A\times C\to C$ 是典范态射。

\begin{proposition}\label{it:1.1.3}
设
\[
0\to A_1\xrightarrow{i_1}A\xrightarrow{p_2}A_2\to 0
\]
是阿贝尔范畴中的短正合序列。则以下条件等价。
\begin{enumerate}[label=(\roman*),leftmargin=2.8em]
  \item 上述短正合序列分裂。
  \item 存在态射 $p_1\colon A\to A_1$，使 $p_1i_1=\id_{A_1}$。
  \item 存在态射 $i_2\colon A_2\to A$，使 $p_2i_2=\id_{A_2}$。
\end{enumerate}
\end{proposition}

\begin{proof}
(i)$\Rightarrow$(ii) 与 (i)$\Rightarrow$(iii) 显然。

(ii)$\Rightarrow$(i) 考虑态射 $\id_A-i_1p_1\colon A\to A$。有
\[
(\id_A-i_1p_1)i_1=i_1-i_1(p_1i_1)=0.
\]
因 $A_2$ 是 $i_1\colon A_1\to A$ 的余核，存在态射 $i_2\colon A_2\to A$ 使 $\id_A-i_1p_1=i_2p_2$。有
\[
p_2i_2p_2=p_2\id_A-p_2i_1p_1=p_2=\id_{A_2}p_2.
\]
因 $p_2$ 是满射，必有 $p_2i_2=\id_{A_2}$。我们的结论于是由 \ref{it:1.1.2}(ii) 推出。

类似地可证 (iii)$\Rightarrow$(i)。
\end{proof}

设 $F\colon\sC\to\sD$ 是阿贝尔范畴之间的共变函子。若对 $\sC$ 中任意对象 $A$ 与 $B$，映射
\[
\Hom(A,B)\to\Hom(F(A),F(B))
\]
是同态，则称 $F$ 是\emph{加性}的。此时有
\[
F(A\oplus B)\cong F(A)\oplus F(B).
\]
事实上，沿用 \ref{it:1.1.2}(i) 的记号并把 $F$ 作用到那里的等式上，得到
\begin{align*}
F(p_1)F(k_1)&=\id_{F(A)},\\
F(p_2)F(k_2)&=\id_{F(B)},\\
F(k_1)F(p_1)+F(k_2)F(p_2)&=\id_{F(A\oplus B)}.
\end{align*}
于是由 \ref{it:1.1.2}(ii)，$(F(A\oplus B),F(k_1),F(k_2))$ 是 $F(A)$ 与 $F(B)$ 的直和。因此若
\[
0\to A\to B\to C\to 0
\]
是分裂短正合序列，则
\[
0\to F(A)\to F(B)\to F(C)\to 0
\]
也是分裂短正合序列。

注意交换群范畴是阿贝尔范畴。我们把以下命题的证明留给读者。
